% =====================================================================
%  Yuji Tachikawa, "A brief review of the 2d/4d correspondences"
%  arXiv:1608.02964v5 (Contribution to "Localization techniques in
%  quantum field theories", eds. Pestun & Zabzine) —— 中文译注本
%
%  编译：LuaLaTeX（本机 xdvipdfmx 不稳定，统一用 lualatex，编译 3 遍）
%  用法：lualatex -interaction=nonstopmode -file-line-error main.tex  （×3）
%        或直接跑 build.ps1（会镜像到 ASCII 目录再编译）
%
%  约定（详见 work/STYLE_GUIDE.md）：
%   * 公式编号一律 \tag{} 复刻原文编号（如 (2.13)），与原文 137 条一一对应；
%   * 图 1--13 为原文矢量图按 500 dpi 截取（figN.png），图注为译者所译；
%   * % [Tutor Note: ...] 为译者批注，不进入排版，汇总见 译注汇总.md；
%   * 章节编号由 \section/\subsection/\subsubsection 自动生成，与原文一致。
% =====================================================================
\documentclass[11pt,a4paper]{article}

\usepackage[UTF8]{ctex}
\usepackage{amsmath,amssymb,amsthm,mathtools}
\usepackage{bm}
\usepackage{graphicx}
\usepackage[a4paper,left=2.4cm,right=2.4cm,top=2.6cm,bottom=2.7cm]{geometry}
\usepackage[dvipsnames]{xcolor}
\usepackage{enumitem}
\usepackage{booktabs}
\usepackage{array}
\usepackage{multirow}
\usepackage{longtable}
\usepackage{multicol}
\usepackage{fancyhdr}
\usepackage{caption}
\usepackage{tcolorbox}
\usepackage{tikz}
\usepackage[hidelinks,unicode]{hyperref}
\hypersetup{pdftitle={2d/4d 对应关系简述（中文译注本）},
  pdfauthor={Yuji Tachikawa; 中译/批注: 本译本},
  hypertexnames=false}

\captionsetup{font=small,labelsep=quad}
\setlength{\emergencystretch}{3em}
\allowdisplaybreaks[1]
\graphicspath{{figures/}{../figures/}}

% ---------- 页眉 ----------
\pagestyle{fancy}
\fancyhf{}
\fancyhead[L]{\small\itshape A brief review of the 2d/4d correspondences}
\fancyhead[R]{\small\thepage}
\renewcommand{\headrulewidth}{0.4pt}

% ---------- 定理类环境（编号一律由原文手工给出） ----------
\theoremstyle{plain}
\newtheorem*{theorem*}{定理}
\newtheorem*{proposition*}{命题}
\newtheorem*{lemma*}{引理}
\theoremstyle{definition}
\newtheorem*{definition*}{定义}
\newtheorem*{remark*}{注}

% 原文的「Exercise x.y.」：粗体标签 + 斜体正文
\newenvironment{exercise}[1]%
  {\par\medskip\noindent\textbf{习题 #1.}\itshape\ }{\par\medskip}

% 原文的灰色要点框
\newtcolorbox{keybox}[1]{%
  colback=black!5, colframe=black!55, colbacktitle=black!75, coltitle=white,
  fonttitle=\bfseries\small, title={#1}, boxrule=0.6pt, arc=2pt,
  left=6pt, right=6pt, top=4pt, bottom=4pt, toptitle=2pt, bottomtitle=2pt,
  fontupper=\small}

% 无标题版要点框
\newtcolorbox{keyboxplain}{%
  colback=black!5, colframe=black!55, boxrule=0.6pt, arc=2pt,
  left=6pt, right=6pt, top=4pt, bottom=4pt, fontupper=\small}

% 原文的粗体行内小标题（如 "Tame punctures."）——独立成段、粗体、不缩进
\newcommand{\para}[1]{\par\medskip\noindent\textbf{#1}\quad\ignorespaces}

% 原文图/表标题（手工编号）。必须用 \captionof*：译者在 \begin{center} 里调用，
% 不是浮动体，用 \caption* 会报 "caption outside float" 且图注文字丢失。
% 开头的 \par\nobreak 把图注与上一行的图像粘住——在 center 环境里图像和图注是两个
% 段落，中间可被分页拆开（图 4、图 12 曾出现「图在上页底、图注跑到下页顶」）。
\newcommand{\tabtitle}[2]{\par\nobreak\captionof*{table}{\small\textbf{表 #1：}#2}}
\newcommand{\figtitle}[2]{\par\nobreak\captionof*{figure}{\small\textbf{图 #1：}#2}}

% 参考文献条目
\newcommand{\refitem}[2]{\par\hangindent=1.5em\hangafter=1\noindent [#1]~#2}

% ---------- 常用记号（只能用这里定义过的宏，不要自造 \newcommand） ----------
\newcommand{\R}{\mathbb{R}}
\newcommand{\C}{\mathbb{C}}
\newcommand{\Z}{\mathbb{Z}}
\newcommand{\N}{\mathbb{N}}
\newcommand{\Q}{\mathbb{Q}}
\newcommand{\E}{\mathbb{E}}
\newcommand{\cA}{\mathcal{A}}
\newcommand{\cB}{\mathcal{B}}
\newcommand{\cC}{\mathcal{C}}
\newcommand{\cD}{\mathcal{D}}
\newcommand{\cE}{\mathcal{E}}
\newcommand{\cF}{\mathcal{F}}
\newcommand{\cG}{\mathcal{G}}
\newcommand{\cH}{\mathcal{H}}
\newcommand{\cI}{\mathcal{I}}
\newcommand{\cJ}{\mathcal{J}}
\newcommand{\cK}{\mathcal{K}}
\newcommand{\cL}{\mathcal{L}}
\newcommand{\cM}{\mathcal{M}}
\newcommand{\cN}{\mathcal{N}}
\newcommand{\cO}{\mathcal{O}}
\newcommand{\cP}{\mathcal{P}}
\newcommand{\cQ}{\mathcal{Q}}
\newcommand{\cR}{\mathcal{R}}
\newcommand{\cS}{\mathcal{S}}
\newcommand{\cT}{\mathcal{T}}
\newcommand{\cU}{\mathcal{U}}
\newcommand{\cV}{\mathcal{V}}
\newcommand{\cW}{\mathcal{W}}
\newcommand{\cX}{\mathcal{X}}
\newcommand{\cY}{\mathcal{Y}}
\newcommand{\cZ}{\mathcal{Z}}
\newcommand{\fg}{\mathfrak{g}}
\newcommand{\fh}{\mathfrak{h}}
\newcommand{\ft}{\mathfrak{t}}
\newcommand{\fu}{\mathfrak{u}}
\newcommand{\fsu}{\mathfrak{su}}
\newcommand{\fso}{\mathfrak{so}}
\newcommand{\fsp}{\mathfrak{sp}}
\newcommand{\SU}{\mathrm{SU}}
\newcommand{\SO}{\mathrm{SO}}
\newcommand{\Sp}{\mathrm{Sp}}
\newcommand{\USp}{\mathrm{USp}}
\newcommand{\U}{\mathrm{U}}
\newcommand{\GL}{\mathrm{GL}}
\newcommand{\dd}{\mathrm{d}}
\newcommand{\ee}{\mathrm{e}}
\newcommand{\ii}{\mathrm{i}}
\newcommand{\Tr}{\operatorname{Tr}}
\newcommand{\tr}{\operatorname{tr}}
\newcommand{\diag}{\operatorname{diag}}
\newcommand{\sgn}{\operatorname{sgn}}
\newcommand{\ord}{\operatorname{ord}}
\newcommand{\vol}{\operatorname{vol}}
\newcommand{\rk}{\operatorname{rk}}
\newcommand{\rank}{\operatorname{rank}}
\newcommand{\Pf}{\operatorname{Pf}}
\newcommand{\Res}{\operatorname{Res}}
\newcommand{\id}{\mathrm{id}}
\renewcommand{\Re}{\operatorname{Re}}
\renewcommand{\Im}{\operatorname{Im}}
\newcommand{\Toda}{\mathrm{Toda}}
\newcommand{\Liouville}{\mathrm{Liouville}}
\newcommand{\half}{\tfrac{1}{2}}
\newcommand{\ie}{\textit{i.e.}}
\newcommand{\eg}{\textit{e.g.}}
\newcommand{\cf}{\textit{cf.}}
% 带上划线的简写（原文的 \bar 系列）
\newcommand{\bQ}{\bar{Q}}
\newcommand{\bq}{\bar{q}}
\newcommand{\bN}{\bar{N}}
\newcommand{\bepsilon}{\bar{\epsilon}}
\newcommand{\btheta}{\bar{\theta}}
\newcommand{\bphi}{\bar{\phi}}
\newcommand{\bpsi}{\bar{\psi}}
\newcommand{\blambda}{\bar{\lambda}}
